% Ad Familiares 7.14 --- headnote
% ad-familiares-07-14, mid-July 53 BC, Rome

\textit{Cicero to C. Trebatius Testa in Gaul, written from Rome
about mid-July 53 BC. The occasion is a verbal message: Chrysippus
Vettius, the freedman of Cyrus --- Cicero's late architect, who
had built or remodelled several of his houses --- has carried oral
greetings from Trebatius, but no letter. Cicero turns the absence
into the joke of the piece: Trebatius has grown too grand to put
pen to paper for ``a man practically of his own household,'' and
the loss to litigation is real --- ``the fewer suitors will be
losing their cases with you as their counsel.'' The threat to
come out and visit before dropping out of memory altogether, and
the renewed dig at Trebatius's evasion of the summer campaigning
(``contrive something, as you did over Britain''), keep the
familiar Trebatius register: affectionate teasing of the young
jurist who has talked himself onto Caesar's staff and is busily
talking himself off the front line.}

\textit{The second paragraph closes with the running pun on the
two kinds of \textit{ius}. Trebatius has been reported on
familiar terms with Caesar --- excellent news --- but Cicero
would rather hear it from Trebatius himself; that, he says,
is what would happen if his friend had ever bothered to master
the law of friendship instead of the law of lawsuits. The
sign-off (``we have been jesting in your way, and not a little
in our own'') acknowledges the manner of the whole piece, and
the warmth that closes it is unforced: Cicero is fond of this
correspondent, and lets the wit drop for the last sentence.}
